%================================================================
% 作業ログ 2026-06-09（test_v3 ゲームモード Processing UI・旋律修正）
%   先週（6/8 週）は台風で休講・提出不要だったため，
%   先週＋今週の 2 週間分をまとめた作業ログ．
%   test_v3 Processing UI 刷新・マスターリセット検知・楽譜旋律修正・
%   IMU モーション表示のジャイロ化まで．
%
%   ビルド:
%     cd work/shiozawa/work-0609/作業ログ0609
%     docker run --rm -v "$(pwd):/workspace" -w /workspace \
%       ghcr.io/paperist/texlive-ja:debian latexmk 作業ログ0609.tex
%================================================================
\documentclass[uplatex,dvipdfmx,a4paper,11pt]{jsarticle}

% --- 余白 ---
\usepackage[top=25mm,bottom=25mm,left=25mm,right=25mm]{geometry}

% --- 表・箇条書き・図・色 ---
\usepackage{array}
\usepackage{tabularx}
\usepackage{booktabs}
\usepackage{enumitem}
\usepackage{graphicx}
\usepackage{url}
\usepackage[table]{xcolor}
\usepackage{amssymb}

% --- 段落間隔（Wordのような見た目に寄せる）---
\setlength{\parindent}{0pt}
\setlength{\parskip}{0.6em}

% --- セクション見出しを「番号. 太字」+ 章末に水平線を入れる ---
\usepackage{titlesec}
\titleformat{\section}
  {\normalfont\Large\bfseries}{\thesection.}{0.6em}{}
\titleformat{\subsection}
  {\normalfont\large\bfseries}{\thesubsection}{0.6em}{}
\titlespacing*{\section}{0pt}{1.2em}{0.6em}
\titlespacing*{\subsection}{0pt}{1.0em}{0.4em}

% --- 章末水平線（手で \sectionrule と書いて区切る）---
\newcommand{\sectionrule}{\par\medskip\noindent\rule{\linewidth}{0.6pt}\par\medskip}

% --- 箇条書きの記号を Word 風（・→○→数字）に揃える ---
\renewcommand{\labelitemi}{\textbullet}
\renewcommand{\labelitemii}{\textopenbullet}
\setlist[itemize]{leftmargin=1.6em,itemsep=0.1em,topsep=0.2em}
\setlist[enumerate]{leftmargin=2.0em,itemsep=0.1em,topsep=0.2em}

% \textopenbullet が無い環境向けのフォールバック
\providecommand{\textopenbullet}{\ensuremath{\circ}}

% --- 進捗ガントチャート用（グリッド表）---
\definecolor{barDone}{gray}{0.55} % 完了
\definecolor{barProg}{gray}{0.78} % 進行中
\definecolor{barPlan}{gray}{0.90} % 予定
\definecolor{nowCol}{gray}{0.30}  % 今週の見出し
\newcolumntype{G}{>{\centering\arraybackslash}p{2.0em}}
\newcommand{\bD}{\cellcolor{barDone}}
\newcommand{\bP}{\cellcolor{barProg}}
\newcommand{\bQ}{\cellcolor{barPlan}}

\begin{document}

%================================================================
% ヘッダ
%================================================================
{\LARGE\bfseries 作業ログ}\par\bigskip

報告者：塩澤 匠生

報告日：2026年6月9日

プロジェクト名：タクトーン~ IMUジェスチャーで奏でるArduinoオーケストラ~ \\

\sectionrule

%================================================================
% 1. 進捗サマリー
%================================================================
\section{進捗サマリー（要約）}

\begin{itemize}
  \item 全体進捗率：60\%（実装フェーズの仕上げ段階．Processing UI は大幅に進捗し，
        firmware 側も旋律修正・復帰高速化・ジャイロ化で改善．実機での統合テストがまだ）
  \item 進捗状況（この 2 週間の変化）：
    \begin{itemize}
      \item test\_v3 Processing UI：メニュー／自由演奏／ゲーム／結果の 4 画面遷移が動作．
            梅澤 UI スタイル適用（淡い青背景＋白パネル），マスターリセット検知，指揮棒軌跡アニメーション，
            2D IMU グラフまで実装
      \item test\_v3 firmware：かえるのうた旋律を 24 拍→32 拍 4 フレーズに修正，
            モーション表示を dynAcc→gyro に変更して安定化，マスターリセット復帰を 5 秒→2 秒に高速化，
            8 分音符のギャップ（0.03 拍）を解消
    \end{itemize}
  \item 今週の主要成果：
    \begin{itemize}
      \item Processing 画面に梅澤 UI スタイル（淡い青背景＋装飾円＋グリッド，白パネル＋影＋青枠，
            モードボタン）を適用し，ウィンドウを 1000$\times$560 に拡大
      \item マスター ESP32 のリセット検知を実装（UI パケット 5 秒→2 秒タイムアウトで待機画面に戻し，
            復帰時は自動追従）
      \item 指揮棒の動きを BPM ベースの振り子アニメーションとして全画面に重ね描画し，
            NOTE 受信時にパルス演出を追加
      \item Conducting 時に IMU 加速度→ジャイロ（角速度$\div$8）の 2D XY プロットを右上に表示
            （80→120 フレームリングバッファ）
      \item かえるのうた旋律を正しい 32 拍 4 フレーズ（ドレミファミレドー／ミファソラソファミー／
            グワッ$\times$4／ドドレレミミファファミレドー）に修正
    \end{itemize}
  \item 重大課題（影響度：中）：
    \begin{itemize}
      \item {}[全 5 ノードでビルドは通るが，実機での統合テスト（全ノード書き込み＋通しプレイ）が未実施．
            特にジャイロ化後のモーション表示のスケール（$\div$8）と UI 中継 30Hz の体感を実機で要確認]
    \end{itemize}
\end{itemize}

\sectionrule

%================================================================
% 2. 進捗ガントチャート（計画と現在地）
%================================================================
\section{進捗ガントチャート（計画と現在地）}

チームのガントチャート（\texttt{meetings/ガントチャーver1ト.xlsx}）のフェーズ区分（110 基本設計／
120 詳細設計／200 実装／300 テスト／400 ストレッチ）に，現在の進み具合を重ねたもの．
「今」の縦位置は本報告時点（6 月第 2 週）．

\medskip
{\footnotesize
\setlength{\tabcolsep}{2pt}
\renewcommand{\arraystretch}{1.2}
\begin{tabular}{@{}p{8.5em}|*{12}{G}@{}}
 & & & & & & & & & \multicolumn{1}{c}{$\downarrow$今} & & & \\
作業項目（フェーズ）
 & \cellcolor{barPlan}\tiny 4/15 & \cellcolor{barPlan}\tiny 4/22 & \cellcolor{barPlan}\tiny 4/29
 & \cellcolor{barPlan}\tiny 5/6 & \cellcolor{barPlan}\tiny 5/13 & \cellcolor{barPlan}\tiny 5/20
 & \cellcolor{barPlan}\tiny 5/27 & \cellcolor{barPlan}\tiny 6/3 & \cellcolor{nowCol}\tiny\textcolor{white}{6/10}
 & \cellcolor{barPlan}\tiny 6/17 & \cellcolor{barPlan}\tiny 6/24 & \cellcolor{barPlan}\tiny 7/1 \\
\midrule
企画・テーマ決定        & \bD & \bD &     &     &     &     &     &     &     &     &     &     \\
110 基本設計            &     &     & \bD & \bD &     &     &     &     &     &     &     &     \\
120 詳細設計            &     &     &     &     & \bD &     &     &     &     &     &     &     \\
\textbf{200 実装}       &     &     &     &     &     &     &     &     &     &     &     &     \\
\quad test\_v2 輪唱・同期 &     &     &     & \bD & \bD & \bD & \bD &     &     &     &     &     \\
\quad test\_v3 ゲーム(firmware) &  &     &     &     &     & \bD & \bD & \bD &     &     &     &     \\
\quad test\_v3 Processing &   &     &     &     &     &     & \bD & \bD & \bP &     &     &     \\
\textbf{300 テスト}     &     &     &     &     &     &     &     &     &     &     &     &     \\
\quad 単体・結合・システム &   &     &     &     &     &     &     &     & \bQ & \bQ & \bQ &     \\
\quad 評価（精度・遅延） &     &     &     &     &     &     &     &     & \bQ & \bQ & \bQ &     \\
400 ストレッチ（強弱ほか） &  &     &     &     &     &     &     & \bQ & \bQ &     &     &     \\
発表会                  &     &     &     &     &     &     &     &     &     &     &     & \multicolumn{1}{c}{$\bigstar$} \\
\bottomrule
\end{tabular}
}

\medskip
凡例：\colorbox{barDone}{~~}~完了 \quad \colorbox{barProg}{~~}~進行中 \quad
\colorbox{barPlan}{~~}~予定 \quad $\bigstar$~発表会（7/1）

\medskip
※ test\_v3 Processing は前回（6/1）の進行中から大幅に進み，4 画面遷移・UI スタイル・
マスターリセット検知・IMU グラフまで実装済み．残りは実機テストでの調整．
※ 300 テストフェーズは 6/10 週から予定通り開始する．

\sectionrule

%================================================================
% 3. 作業ログ（詳細トラッキング）
%================================================================
\section{作業ログ（詳細トラッキング）}

\begin{tabularx}{\linewidth}{@{}p{8.5em}p{5em}p{5em}p{4em}p{3em}X@{}}
\toprule
作業項目 & 開始日 & 予定完了日 & 状態 & 進捗 & コメント \\
\midrule
Processing UI 刷新 & 2026-06-07 & 2026-06-07 & 完了 & 100\% & 梅澤 UI スタイル適用（淡い青背景＋白パネル＋青枠），ウィンドウ 1000$\times$560 \\
マスターリセット検知 & 2026-06-07 & 2026-06-07 & 完了 & 100\% & UI パケット 5 秒→2 秒タイムアウトで待機画面に遷移，復帰時自動追従 \\
楽譜旋律修正 & 2026-06-07 & 2026-06-08 & 完了 & 100\% & 24 拍→32 拍 4 フレーズに拡張，正しいかえるのうた旋律に修正 \\
8 分音符修正 & 2026-06-07 & 2026-06-08 & 完了 & 100\% & durationQ8/subDurationQ8 を 120→128 に修正（0.03 拍のギャップ解消） \\
指揮棒軌跡アニメーション & 2026-06-07 & 2026-06-07 & 完了 & 100\% & BPM ベースの振り子を全画面に重ね描画，NOTE 受信時パルス演出 \\
メニュー横配置 & 2026-06-07 & 2026-06-07 & 完了 & 100\% & 縦→横並びに変更（IMU 左右振りと直感的に対応） \\
2D IMU グラフ & 2026-06-08 & 2026-06-08 & 完了 & 100\% & ジャイロ（角速度$\div$8）の XY プロットを右上に表示，120 フレームリングバッファ \\
実機テスト & 2026-06-10 & — & これから & — & 全ノード書き込み，統合テスト，ジャイロスケール・UI 中継の体感確認 \\
\bottomrule
\end{tabularx}

\sectionrule

%================================================================
% 4. 課題管理
%================================================================
\section{課題管理}

\begin{tabularx}{\linewidth}{@{}p{8em}Xp{2.5em}p{2.5em}Xp{3em}@{}}
\toprule
課題 & 発生原因 & 影響度 & 優先度 & 対策 & 状態 \\
\midrule
IMU モーション表示のスケール & dynAcc は重力成分が混入して振りの可視化に不向きだった & 中 & 高 & ジャイロ（角速度$\div$8）に変更して安定化．スケール値は実機で要調整 & 対応中 \\
UDP マルチキャストのパケロス & マルチキャストは ACK・再送が無く，取りこぼしが音抜けになる & 中 & 中 & 連送 4 発の時間分散＋WiFi 復帰時の再 join で対策済み．実機テストで効果確認する & 継続 \\
5 台構成（金管 4＋ドラム）の輪唱位相 & 等間隔で 4 声以上重ねると周期が一巡せず位相が破綻する & 低 & 低 & 不等間隔の位相か曲の周期延長＝編曲（音楽判断）が必要 & 保留 \\
\bottomrule
\end{tabularx}

\medskip
※影響度＝プロジェクトへのインパクト，優先度＝対応の緊急性

\sectionrule

%================================================================
% 5. AI利用状況と検証
%================================================================
\section{AI利用状況と検証(再現性重視)}

\subsection{利用内容}

\begin{itemize}
  \item 利用目的：Processing UI 設計のレビュー，楽譜データ修正，IMU データの可視化方式の検討
  \item 入力（プロンプト要旨）：
    \begin{itemize}
      \item {}[かえるのうたの旋律を正しい 32 拍 4 フレーズに修正する方法]
      \item {}[dynAcc（重力混入）に代わる，指揮棒モーションの安定した可視化方式]
      \item {}[マスター ESP32 のリセット検知と復帰時の画面遷移ロジック]
    \end{itemize}
  \item 出力概要：
    \begin{itemize}
      \item {}[24 拍 3 フレーズ→32 拍 4 フレーズへの拡張と各フレーズの MIDI 音程の具体案]
      \item {}[ジャイロ（角速度$\div$8）を int8 で送り，Processing で 2D XY プロットする案]
      \item {}[UI パケットのタイムアウト（5 秒→2 秒）で待機画面に戻し，復帰時に自動追従する案]
    \end{itemize}
\end{itemize}

\sectionrule

\subsection{検証方法・ファクトチェック}

\begin{itemize}
  \item 検証手法：
    \begin{itemize}
      \item {}[全 5 ノードを \texttt{pio run} でビルドし，コンパイル成功を確認（全コミットで SUCCESS）]
      \item {}[GAME\_LENGTH\_BEATS を 24→32 に変更し，score\_data.cpp のフレーズ数と整合していることをコードで確認]
    \end{itemize}
  \item 最終判断：
    \begin{itemize}
      \item 採用．ジャイロ化と旋律修正は机上では正しいが，実機でのモーション表示の体感と
            8 分音符の聞こえ方は実機テストで最終確認する
    \end{itemize}
\end{itemize}

\sectionrule

%================================================================
% 6. 次回計画
%================================================================
\section{次回計画}

\begin{itemize}
  \item 目標：
    \begin{itemize}
      \item {}[全ノードへ firmware を書き込み，実機での統合テストを行う．
            4 画面遷移（メニュー→自由演奏→ゲーム→結果）が通しで動くことを確認する]
      \item {}[ジャイロスケール（$\div$8）と UI 中継 30Hz の体感を実機で確認し，必要に応じて調整する]
      \item {}[テストフェーズ（300）に移行し，精度・遅延の評価を開始する]
      \item {}[発表準備（スライド・デモシナリオ）を開始する]
    \end{itemize}
  \item 達成基準：
    \begin{itemize}
      \item {}[実機で自由演奏モードが test\_v2 と同等の安定性で鳴る]
      \item {}[実機でメニュー→ゲーム→結果の遷移が破綻なく動き，採点結果が表示される]
    \end{itemize}
  \item 期限：
    \begin{itemize}
      \item {}[6/17 までに実機テスト・調整を完了し，テストフェーズに本格移行]
    \end{itemize}
\end{itemize}

\sectionrule

%================================================================
% 7. その他
%================================================================
\section{その他}

\begin{itemize}
  \item {}[先週（6/8 週）は台風により休講・作業ログ提出不要だったため，先週＋今週の 2 週間分をまとめて報告している]
  \item {}[ゲームモードは追加機能（ストレッチ）であり，遅れても自由演奏のみで発表は成立する位置づけ]
  \item {}[実機未テストのファームウェアへの机上修正は同期挙動を逆転させるリスクがあるため，必ず実機書き込みでの確認を経て取り込む]
\end{itemize}

\sectionrule

%================================================================
% 8. 付録
%================================================================
\section{付録}

\begin{itemize}
  \item 添付資料：
    \begin{itemize}
      \item ソースコード：\texttt{firmware/test\_v3/}（共通ライブラリ＋node\_01〜node\_04），\texttt{pc\_app/test\_v3/}
      \item GitHub リポジトリ：\url{https://github.com/takushio2525/2026-hackathon-team23}
      \item 関連コミット：\texttt{4320239}，\texttt{20a0451}，\texttt{f85e63d}，\texttt{cfea5ea}
    \end{itemize}
  \item 添付範囲：
    \begin{itemize}
      \item 差分（test\_v3 Processing UI 刷新・旋律修正・ジャイロ化の直近 2 週間分のスナップショット）
    \end{itemize}
\end{itemize}

\sectionrule

\end{document}
